\documentclass[12pt, twoside]{article}
%\documentclass[12pt, twoside]{article}
\usepackage[letterpaper, margin=1in, headsep=0.2in]{geometry}
\setlength{\headheight}{0.6in}
%\usepackage[english]{babel}
\usepackage[utf8]{inputenc}
\usepackage{microtype}
\usepackage{amsmath}
\usepackage{amssymb}
%\usepackage{amsfonts}
\usepackage[nomessages]{fp} %\FPeval{\var-name}{2*sin(pi/6)}
\usepackage{siunitx} %units in math. eg 20\milli\meter
\usepackage{yhmath} % for arcs, overparenth command
\usepackage{tikz} %graphics
\usetikzlibrary{quotes, angles, arrows, arrows.meta}
\usepackage{pgfplots}
\usepgfplotslibrary{fillbetween}
\pgfplotsset{compat=1.16}
\usepackage{graphicx} %consider setting \graphicspath{{images/}}
\usepackage{parskip} %no paragraph indent
\usepackage{enumitem}
\usepackage{multicol}
\usepackage{venndiagram}

\usepackage{fancyhdr}
\pagestyle{fancy}
\fancyhf{}
\renewcommand{\headrulewidth}{0pt} % disable the underline of the header
\raggedbottom
\hfuzz=2mm %suppresses overfull box warnings

\usepackage{hyperref}
\usepackage{eurosym}

\fancyhead[LO]{La Scuola d'Italia / Dr.\ Huson / IB Math 11 \\* 14 October 2026}
\fancyhead[RO]{\fbox{\begin{minipage}{6cm}\footnotesize Nickname:\\[0.5cm]\end{minipage}}}
\fancyhead[LE]{\fbox{\begin{minipage}{6cm}\footnotesize Nickname:\\[0.5cm]\end{minipage}}}
\fancyfoot[C]{\thepage}

\begin{document}
\subsubsection*{1.12 Homework: Review of Sequences and Series}

\begin{enumerate}[itemsep=0.15cm]

%--- Problem 1: Arithmetic sequence; d, u_n, u_25, which term is 143 ---
%--- Source: Original (freshly authored, AA SL 1.2); review of 1-1CW, 1-3CW ---
\item[1.]
Consider the arithmetic sequence
$$7,\; 11,\; 15,\; 19,\; \ldots$$

\begin{enumerate}[itemsep=0.1cm]
  \item Write down the common difference $d$. \vfill
  \item Find an expression for $u_{n}$, the $n^{\text{th}}$ term. \vfill
  \item Find $u_{25}$. \vfill
  \item Find which term of the sequence is equal to $143$. \vfill
\end{enumerate}

%--- Problem 2: Arithmetic sequence from two terms; u_1, d, u_20, S_20 ---
%--- Source: Original (freshly authored, AA SL 1.2); review of 1-2CW, 1-5Q ---
\item[2.]
In an arithmetic sequence, $u_{3} = 14$ and $u_{8} = 39$.

\begin{enumerate}[itemsep=0.1cm]
  \item Find the common difference $d$ and the first term $u_{1}$. \vfill
  \item Find $u_{20}$. \vfill
  \item Find $S_{20}$, the sum of the first $20$ terms. \vfill
\end{enumerate}

\newpage

%--- Problem 3: Arithmetic series with a given last term; number of terms, sum ---
%--- Source: Original (freshly authored, AA SL 1.2); review of 1-3CW_Series ---
\item[3.]
Consider the arithmetic series
$$5 + 9 + 13 + 17 + \cdots + 201.$$

\begin{enumerate}[itemsep=0.1cm]
  \item Find the number of terms in the series. \vfill
  \item Find the sum of the series. \vfill
\end{enumerate}

%--- Problem 4: Theatre seats; u_25, S_25, least number of rows for 600 seats ---
%--- Source: Original (freshly authored, AA SL 1.2); review of 1-5Q ---
\item[4.]
A theatre has $25$ rows of seats. There are $20$ seats in the first row,
and each row has $2$ more seats than the row in front of it.

\begin{enumerate}[itemsep=0.1cm]
  \item Find the number of seats in the last row. \vfill
  \item Find the total number of seats in the theatre. \vfill
  \item A school group needs at least $600$ seats. The group fills the
  rows in order, starting from the first row. Find the least number of
  rows the group needs. \vfill
\end{enumerate}

\newpage

%--- Problem 5: Geometric sequence; r, u_n, u_10, first term over 10 000 ---
%--- Source: Original (freshly authored, AA SL 1.3); review of 1-6CW, 1-9PQ ---
\item[5.]
Consider the geometric sequence
$$3,\; 6,\; 12,\; 24,\; \ldots$$

\begin{enumerate}[itemsep=0.1cm]
  \item Write down the common ratio $r$. \vfill
  \item Find an expression for $u_{n}$, the $n^{\text{th}}$ term. \vfill
  \item Find $u_{10}$. \vfill
  \item Find the first term of the sequence that is greater than
  $10\,000$, and state which term it is. \vfill
\end{enumerate}

%--- Problem 6: Geometric sequence from two terms; r, u_1, u_7, S_6 ---
%--- Source: Original (freshly authored, AA SL 1.3); review of 1-9PQ, 1-11CW ---
\item[6.]
In a geometric sequence, $u_{2} = 12$ and $u_{5} = 324$.

\begin{enumerate}[itemsep=0.1cm]
  \item Find the common ratio $r$ and the first term $u_{1}$. \vfill
  \item Find $u_{7}$. \vfill
  \item Find $S_{6}$, the sum of the first six terms. \vfill
\end{enumerate}

\newpage

%--- Problem 7: Training distance growing 10% per week; r, week-8 distance, 8-week total ---
%--- Source: Original (freshly authored, AA SL 1.3); review of 1-6CW, 1-11CW ---
\item[7.]
Giulia is training for a race. She runs $5$ km in the first week, and
each week she increases her distance by $10\%$.

\begin{enumerate}[itemsep=0.1cm]
  \item Write down the common ratio of the weekly distances. \vfill
  \item Find the distance she runs in the eighth week. \vfill
  \item Find the total distance she runs in the first eight weeks. \vfill
\end{enumerate}

%--- Problem 8: Compound interest, quarterly vs annual; interest earned; years to exceed 8000 ---
%--- Source: Original (freshly authored, AA SL 1.4); review of 1-7CW, 1-8CW ---
\item[8.]
Marco invests \euro$5000$ in a savings account that pays $4\%$ per year,
compounded \textbf{quarterly}.

\begin{enumerate}[itemsep=0.1cm]
  \item Find the value of the investment after $6$ years. \vfill
  \item Find the interest Marco earns in these $6$ years. \vfill
  \item Another account pays $4\%$ per year, compounded
  \textbf{annually}. Find the value of \euro$5000$ invested in this
  account after $6$ years. \vfill
  \item For the quarterly account, find the number of complete years
  after which the value of Marco's investment first exceeds
  \euro$8000$. \vfill
\end{enumerate}

\end{enumerate}
\end{document}
